% chapters/10/section-01.tex
\section{Peta Berpetak sebagai Graf}
\label{sec:peta-graf}

Bagian~2.5 telah menjelaskan bahwa pencarian jalur adalah persoalan pada graf.
Bagian ini mengubah peta permainan menjadi graf tersebut.

\subsection{Dari Gambar Menjadi Simpul dan Sisi}

Sebuah peta yang tampak sebagai gambar sesungguhnya tersusun dari petak-petak.
Setiap petak kosong dapat dilalui, sedangkan petak dinding tidak. Perubahan
dari gambar menjadi graf dilakukan dengan tiga aturan sederhana.

\begin{enumerate}
  \item \textbf{Setiap petak kosong menjadi satu simpul}, dinyatakan oleh
        pasangan angka, yaitu kolom dan barisnya.
  \item \textbf{Setiap pasangan petak yang bersebelahan menjadi satu sisi}.
        Pada bab ini hanya empat arah yang dipakai, yaitu atas, bawah, kiri,
        dan kanan, sehingga selalu ada paling banyak empat sisi pada setiap
        simpul.
  \item \textbf{Biaya setiap sisi bernilai satu}, sebab memindahkan benda satu
        petak ke petak sebelahnya selalu sama jauhnya. Apabila kelak Anda
        menghendaki medan yang berbeda-beda, misalnya pasir yang memperlambat,
        biaya sisi itu dapat dibedakan.
\end{enumerate}

Ketiga aturan itu diperlihatkan pada panel kiri
Gambar~\ref{fig:peta-heuristik}.

\begin{figure}[htbp]
\centering
% chapters/10/figures/peta-dan-heuristik.tikz
% Dua hal yang perlu dipahami sebelum menulis pencarian: peta berpetak sebagai
% graf, dan dua cara mengukur perkiraan jarak ke tujuan.
\begin{tikzpicture}[
  sel/.style={rectangle, draw=ubLine, thick, minimum size=7mm, inner sep=0pt,
              font=\scriptsize},
  dinding/.style={sel, fill=ubGray},
  terbuka/.style={sel, fill=white},
  awal/.style={sel, fill=ubGreenBg, draw=ubGreen},
  tujuan/.style={sel, fill=ubRedBg, draw=ubRed},
  panah/.style={-{Stealth[length=2mm]}, thick, draw=ubAccent},
  judul/.style={font=\footnotesize\bfseries, ubPrimary, align=center},
  label/.style={font=\scriptsize, align=left}
]

% ---------- Panel A: peta berpetak sebagai graf ----------
\begin{scope}
  \node[judul, anchor=south] at (1.5,0.35) {peta berpetak sebagai graf};

  \foreach \x in {0,...,4} {
    \foreach \y in {0,...,3} {
      \node[terbuka] at (\x*0.75,-\y*0.75) {};
    }
  }
  \node[dinding] at (1*0.75,0) {};
  \node[dinding] at (3*0.75,-1*0.75) {};
  \node[dinding] at (1*0.75,-2*0.75) {};
  \node[dinding] at (2*0.75,-3*0.75) {};
  \node[awal] at (0,-3*0.75) {S};
  \node[tujuan] at (4*0.75,0) {G};

  % Empat tetangga dari satu petak.
  \draw[panah] (2*0.75,-1*0.75) -- (2*0.75,-1.5*0.75);
  \draw[panah] (2*0.75,-1*0.75) -- (2*0.75,-0.5*0.75);
  \draw[panah] (2*0.75,-1*0.75) -- (2.5*0.75,-1*0.75);
  \draw[panah] (2*0.75,-1*0.75) -- (1.5*0.75,-1*0.75);

  \node[label, anchor=north west, text width=52mm] at (-0.2,-2.55)
    {Setiap petak kosong menjadi satu simpul, dan hubungan dengan petak
     bersebelahan menjadi sisi graf. Empat arah dipakai, sehingga biaya
     perpindahan selalu satu.};
\end{scope}

% ---------- Panel B: dua ukuran perkiraan jarak ----------
\begin{scope}[xshift=7.4cm]
  \node[judul, anchor=south] at (1.5,0.35) {dua ukuran perkiraan jarak};

  \foreach \x in {0,...,4} {
    \foreach \y in {0,...,3} {
      \node[terbuka] at (\x*0.75,-\y*0.75) {};
    }
  }
  \node[awal] at (0,-3*0.75) {S};
  \node[tujuan] at (4*0.75,0) {G};

  % Jalur tangga (Manhattan) dan garis lurus (Euclidean).
  \draw[ubAccent, thick] (0,-3*0.75) -- (4*0.75,-3*0.75) -- (4*0.75,0);
  \draw[ubRed, thick, dashed] (0,-3*0.75) -- (4*0.75,0);

  \node[label, ubAccent, anchor=north west] at (-0.2,-2.55)
    {Manhattan: $h = |\Delta x| + |\Delta y|$};
  \node[label, ubRed, anchor=north west] at (-0.2,-3.05)
    {Euclidean: $h = \sqrt{\Delta x^2 + \Delta y^2}$};
  \node[label, anchor=north west, text width=52mm] at (-0.2,-3.75)
    {Keduanya tidak melebihkan jarak sebenarnya, sehingga keduanya
     menghasilkan jalur terpendek. Manhattan lebih dekat kepada jarak
     sebenarnya, sehingga lebih sedikit petak diperiksa.};
\end{scope}
\end{tikzpicture}

\caption{Peta berpetak sebagai graf dan dua ukuran perkiraan jarak}
\label{fig:peta-heuristik}
\end{figure}

\subsection{Peta Ditulis sebagai Gambar Huruf}

Peta disimpan sebagai rangkaian baris huruf, sehingga bentuknya masih dapat
dikenali ketika dibaca di dalam kode.

\begin{kodekeluaran}
PETA = (
    "################################",
    "#S...#.........#........#......#",
    "#....#.........#........#......#",
    ...
)

CELL = 28        # sisi setiap petak, dalam piksel
\end{kodekeluaran}

Huruf \code{#} menyatakan dinding, titik menyatakan petak kosong, \code{S}
menyatakan titik awal, dan \code{G} menyatakan titik tujuan. Cara penulisan itu
dipilih karena tiga alasan: petanya mudah disunting, bentuknya mudah
diperiksa dengan mata, dan kesalahan penulisan baris langsung terlihat.

Program bab ini memeriksa petanya sebelum dipakai. Pemeriksaan itu meliputi
panjang seluruh baris yang harus sama, serta keberadaan tepat satu titik awal
dan satu titik tujuan. Pemeriksaan semacam itu murah, dan menghindarkan Anda
dari kesulitan mencari sebab kegagalan program yang sebenarnya berasal dari
peta yang salah tulis.

\subsection{Penentuan Petak Tetangga}

Seluruh hubungan antar-simpul diringkas pada satu fungsi yang mengembalikan
petak-petak yang bersebelahan dengan sebuah petak.

\begin{kodekeluaran}
def tetangga(peta_lebar, peta_tinggi, dinding, petak):
    kolom, baris = petak
    hasil = []
    for geser_kolom, geser_baris in ((1, 0), (-1, 0), (0, 1), (0, -1)):
        baru = (kolom + geser_kolom, baris + geser_baris)
        if 0 <= baru[0] < peta_lebar and 0 <= baru[1] < peta_tinggi:
            if baru not in dinding:
                hasil.append(baru)
    return hasil
\end{kodekeluaran}

Fungsi itu menyaring dua hal sekaligus: tetangga yang
keluar dari batas peta, dan tetangga yang berupa dinding. Dengan begitu,
pencarian tidak perlu memikirkan dinding lagi; ia cukup memanggil fungsi itu
dan menerima petak-petak yang memang dapat dilalui.

\begin{catatabox}
Seluruh peta dibaca ulang setiap kali pencarian
dijalankan pada program bab ini. Cara itu memadai untuk keperluan pembelajaran,
tetapi boros apabila pencarian dilakukan berulang kali. Pada Bab~XI, peta akan
dibaca sekali saja lalu disimpan, sebab pencarian jalur di sana dijalankan
berulang kali untuk beberapa NPC.
\end{catatabox}
